\documentclass[11pt,a4paper]{article}

\def\ncourse {A First Course in Functional Analysis}
\def\ntype {problems}

\makeatletter

% packages
\usepackage{amssymb,amsfonts,amsmath,calc,tikz,pgfplots,geometry,mathtools}
\usepackage{color}   % May be necessary if you want to color links
\usepackage[svgnames]{xcolor}
\usepackage[hidelinks]{hyperref}
\usepackage{forest}
\usepackage{commath} % ew
\usepackage{esvect} % \vv makes vectors
\usepackage{centernot} % for the comparison
\usepackage{esdiff}
\usepackage{esint}
\usepackage{amsthm}
\usepackage{fancyhdr}
\usepackage{bm}
\usepackage{witharrows}
\usepackage{bookmark}
\usepackage{tikz-cd}
\usepackage{quiver}
\usepackage{bbm}
\usepackage{textcomp}
\usepackage{float}
\usepackage{gensymb}
\usepackage{cleveref}
\usepackage{environ}
\usepackage{comment}
\usepackage{cancel}
\usepackage{xparse}

% Inevitable cs
\usepackage{listings}
\usepackage[]{algorithm2e}

% tikz libraries
\usetikzlibrary{positioning}
\usetikzlibrary{matrix}
\usetikzlibrary{arrows}
\usetikzlibrary{bending}
\usetikzlibrary{arrows.meta}
\usetikzlibrary{decorations.markings}
\usetikzlibrary{decorations.pathreplacing}
\usetikzlibrary{decorations.markings}
\usetikzlibrary{patterns}
\usetikzlibrary{backgrounds}

% Page style setup
\pagestyle{fancy}
\fancyhfoffset{0pt}
\renewcommand{\headwidth}{\textwidth}
\geometry{margin=1in}
\pgfplotsset{compat=1.18}
\setlength{\headheight}{14.6pt}
\addtolength{\topmargin}{-1.6pt}
\hypersetup{
    colorlinks=false,
    linktoc=section,
    linkcolor=black,
}

%% maketitle setup
\ifx \nauthor\undefined
  \def\nauthor{Yeheli Fomberg}
\else
\fi

\ifx \nemail\undefined
  \def\nemail{\href{mailto:yp@campus.technion.ac.il}
  {\texttt{<yp@campus.technion.ac.il>}}}
\else
\fi

\ifx \ncoursehead \undefined
  \def\ncoursehead{\ncourse}
\fi

\lhead{\emph{\nouppercase{\leftmark}}}
\ifx \nextra \undefined
  \rhead{
    \ifnum\thepage=1
    \else
      \ncoursehead
    \fi}
\else
  \rhead{
    \ifnum\thepage=1
    \else
      \ncoursehead \ (\nextra)
    \fi}
\fi

\def\notes{notes}
\def\problems{problems}

\ifx \ntype \undefined
  \def\ntype{notes}
\else
\fi

\ifx \ntype \notes
  \let\@real@maketitle\maketitle
  \renewcommand{\maketitle}{\@real@maketitle\begin{center}
  \begin{minipage}[c]{0.9\textwidth}\centering\footnotesize
  These notes are not endorsed by the lecturers.
  I have revised them outside lectures to incorporate supplementary
  explanations, clarifications, and material for fun.
  While I have strived for accuracy, any errors or misinterpretations 
  are most likely mine.
  \end{minipage}\end{center}}
\else
  \ifx \ntype \problems
    \let\@real@maketitle\maketitle
    \renewcommand{\maketitle}{\@real@maketitle\begin{center}
    \begin{minipage}[c]{0.9\textwidth}\centering\footnotesize
    These problems are mostly based on the problem sets I was given
    when taking the course.
    
    While I have strived for accuracy and completeness,
    some of the problems would not have been here without the help of
    friends.
    \end{minipage}\end{center}}
  \else
  \fi
\fi


% theorem environments
\theoremstyle{definition}
\newtheorem{definition}{Definition}[section]
\newtheorem{remark}{Remark}[section]
\newtheorem{example}{Example}[section]
\newtheorem{exercise}{Exercise}[section]
\newtheorem{paradox}{Paradox}[section]
\newtheorem{entry}{Entry}[section]
\newtheorem*{solution}{Solution}
\newtheorem*{question}{Question}
\newtheorem*{answer}{Answer}
\newtheorem*{problem}{Problem}
\theoremstyle{plain}
\newtheorem{theorem}{Theorem}[section]
\newtheorem{proposition}[theorem]{Proposition}
\newtheorem{lemma}[theorem]{Lemma}
\newtheorem{corollary}[theorem]{Corollary}
\newenvironment{proofof}[1]{%
    \renewcommand{\proofname}{Proof of #1}%
    \begin{proof}
}{%
    \end{proof}
}
% tikz customization
\pgfarrowsdeclarecombine{twolatex'}{twolatex'}{latex'}{latex'}{latex'}{latex'}
\tikzset{->/.style = {decoration={markings,
                                  mark=at position 0.999
                                  with {\arrow[scale=2]{latex'}}},
                      postaction={decorate}}}
\tikzset{<-/.style = {decoration={markings,
                                  mark=at position 0.001 with {\arrowreversed[scale=2]{latex'}}},
                      postaction={decorate}}}
\tikzset{-->/.style = {decoration={markings,
                                  mark=at position 0.999
                                  with {\arrow[scale=2]{latex'[sep=0pt -2.5]}}},
                      postaction={decorate}}}
\tikzset{<--/.style = {decoration={markings,
                                  mark=at position 0.001
                                  with {\arrowreversed[scale=2]{latex'[sep=0pt -2.5]}}},
                      postaction={decorate}}}
\tikzset{<->/.style = {decoration={markings,
                                  mark=at position 0.001 with {\arrowreversed[scale=2]{latex'}},
                                  mark=at position 0.999 with {\arrow[scale=2]{latex'}},},
                      postaction={decorate}}}
\tikzset{->-/.style = {decoration={markings,
    mark=at position 0.50+0.225/#1 with {\arrow[scale=2]{latex'}}},
                      postaction={decorate}}}
\tikzset{-<-/.style = {decoration={markings,
                                   mark=at position 0.50-0.225/#1 with {\arrowreversed[scale=2]{latex'}}},
                       postaction={decorate}}}
\tikzset{->>/.style = {decoration={markings,
                                  mark=at position 1 with {\arrow[scale=2]{latex'}}},
                      postaction={decorate}}}
\tikzset{<<-/.style = {decoration={markings,
                                  mark=at position 0 with {\arrowreversed[scale=2]{twolatex'}}},
                      postaction={decorate}}}
\tikzset{<<->>/.style = {decoration={markings,
                                   mark=at position 0 with {\arrowreversed[scale=2]{twolatex'}},
                                   mark=at position 1 with {\arrow[scale=2]{twolatex'}}},
                       postaction={decorate}}}
\tikzset{->>-/.style = {decoration={markings,
                                   mark=at position 0.50+0.450/#1 with {\arrow[scale=2]{twolatex'}}},
                       postaction={decorate}}}
\tikzset{-<<-/.style = {decoration={markings,
                                   mark=at position 0.50-0.450/#1 with {\arrowreversed[scale=2]{twolatex'}}},
                       postaction={decorate}}}

\pgfarrowsdeclare{biggertip}{biggertip}{%
  \setlength{\arrowsize}{1pt}
  \addtolength{\arrowsize}{.1\pgflinewidth}
  \pgfarrowsrightextend{0}
  \pgfarrowsleftextend{-5\arrowsize}
}{%
  \setlength{\arrowsize}{1pt}
  \addtolength{\arrowsize}{.1\pgflinewidth}
  \pgfpathmoveto{\pgfpoint{-5\arrowsize}{4\arrowsize}}
  \pgfpathlineto{\pgfpointorigin}
  \pgfpathlineto{\pgfpoint{-5\arrowsize}{-4\arrowsize}}
  \pgfusepathqstroke
}
\tikzset{
	EdgeStyle/.style = {>=biggertip}
}

\tikzset{circ/.style = {fill, circle, inner sep = 0, minimum size = 3}}
\tikzset{scirc/.style = {fill, circle, inner sep = 0, minimum size = 1.5}}
\tikzset{mstate/.style={circle, draw, black, text=black, minimum width=0.7cm}}

\tikzset{eqpic/.style={baseline={([yshift=-.5ex]current bounding box.center)}}}

\definecolor{mblue}{rgb}{0.2, 0.3, 0.8}
\definecolor{morange}{rgb}{1, 0.5, 0}
\definecolor{mgreen}{rgb}{0, 0.4, 0.2}
\definecolor{mred}{rgb}{0.5, 0, 0}

% topology
\DeclareMathOperator{\dfr}{dfr} % deformation retract
\newcommand{\Cells}{\text{Cells}}
\newcommand{\RP}{\mathbb{RP}}
\newcommand{\CP}{\mathbb{CP}}

% order theory
\newcommand{\join}{\vee}
\newcommand{\meet}{\wedge}

% algebraic geometry
\DeclareMathOperator{\Rad}{Rad} % Radical of modules
\DeclareMathOperator{\Spec}{Spec}
% \renewcommand{\div}{\mathrm{div}} % divisors
\DeclareMathOperator{\Mer}{Mer} % Meromorphic functions
\DeclareMathOperator{\Exc}{Exc} % Exceptional divisor
\DeclareMathOperator{\Td}{Td} % Todd classes

% algebra

\DeclareMathOperator{\trdeg}{tr.deg.}
\DeclareMathOperator{\odd}{odd}
\DeclareMathOperator{\even}{even}
\DeclareMathOperator{\lcm}{lcm}
\DeclareMathOperator{\ord}{ord}
\DeclareMathOperator{\codim}{codim}
\DeclareMathOperator{\Ann}{Ann}
\DeclareMathOperator{\Ext}{Ext}
\DeclareMathOperator{\Sym}{Sym}
\DeclareMathOperator{\Out}{Out}
\DeclareMathOperator{\Aut}{Aut}
\DeclareMathOperator{\End}{End}
\DeclareMathOperator{\Inn}{Inn}
\DeclareMathOperator{\Mat}{Mat}
\DeclareMathOperator{\Fix}{Fix}
\DeclareMathOperator{\std}{std}
\DeclareMathOperator{\sgn}{sgn}
\DeclareMathOperator{\adj}{adj}
\DeclareMathOperator{\id}{id}
\DeclareMathOperator{\op}{op}
\DeclareMathOperator{\tr}{tr}
\DeclareMathOperator{\diag}{diag}
\DeclareMathOperator{\rank}{rank}
\DeclareMathOperator{\rk}{rk}
\DeclareMathOperator{\coker}{coker}
\DeclareMathOperator{\Mult}{Mult} % multiplier algebra
\DeclareMathOperator{\Frac}{Frac} % fraction field
\DeclareMathOperator{\Rat}{Rat}
\DeclareMathOperator{\Gal}{Gal}
\newcommand{\ioplus}{\overset{\text{int}}{\oplus}} % interior direct sum
\newcommand{\GG}{\mathcal G} % Galois correspondence
\newcommand{\FF}{\mathcal F} % Galois correspondence
\newcommand{\cupp}{\smile} % cup product
\newcommand{\capp}{\frown} % cap product
\newcommand{\actson}{\curvearrowright}
\newcommand{\bra}{\left\langle}
\newcommand{\ket}{\right\rangle}
\newcommand{\idealin}{\triangleleft\,}
\newcommand{\normalin}{\triangleleft\,}
\newcommand{\ip}[2]{\left\langle #1, #2 \right\rangle}
\newcommand{\bigslant}[2]
{{\raisebox{.2em}{$#1$}\left/\raisebox{-.2em}{$#2$}\right.}}
\newcommand{\properideal}{%
  \mathrel{\ooalign{$\lneq$\cr\raise.22ex\hbox{$\lhd$}\cr}}}

% categories
\newcommand{\cat}[1]{\mathbf{#1}}
\newcommand{\Ch}{\cat{Ch}} % Chain complexes   
\newcommand{\Set}{\cat{Set}}
\newcommand{\Grp}{\cat{Grp}}
\newcommand{\Ab}{\cat{Ab}}
\newcommand{\Ring}{\cat{Ring}}
\newcommand{\CRing}{\cat{CRing}}
\newcommand{\Top}{\cat{Top}}
\newcommand{\Vect}{\cat{Vect}}
\newcommand{\cmod}{\cat{mod}}

%% matrix groups
\newcommand{\GL}{\mathrm{GL}}
\newcommand{\Or}{\mathrm{O}}
\newcommand{\PGL}{\mathrm{PGL}}
\newcommand{\PSL}{\mathrm{PSL}}
\newcommand{\PSO}{\mathrm{PSO}}
\newcommand{\PSU}{\mathrm{PSU}}
\newcommand{\SL}{\mathrm{SL}}
\newcommand{\msl}{\mathrm{sl}}
\newcommand{\SO}{\mathrm{SO}}
\newcommand{\Spin}{\mathrm{Spin}}
\newcommand{\Sp}{\mathrm{Sp}}
\newcommand{\SU}{\mathrm{SU}}
\newcommand{\U}{\mathrm{U}}
\let\char\relax
\newcommand{\char}{\text{char}}

% analysis
\renewcommand{\Re}{\mathrm{Re}}
\renewcommand{\Im}{\mathrm{Im}}
\NewDocumentCommand{\dx}{o}{
  \IfNoValueTF{#1}
    {\dif x}                  % if no optional argument
    {\dif^{#1} \!x}         % if optional argument
}
\NewDocumentCommand{\dy}{o}{
  \IfNoValueTF{#1}
    {\dif y}                  % if no optional argument
    {\dif^{#1} \!y}         % if optional argument
}
\newcommand{\dt}{\dif t}
\newcommand{\du}{\dif u}
\newcommand{\dv}{\dif v}
\newcommand{\dz}{\dif z}
\newcommand{\ds}{\dif s}
\newcommand{\dw}{\dif w}
\newcommand{\dr}{\dif r}
\newcommand{\dm}{\dif m}
\newcommand{\df}{\dif f}
\newcommand{\dV}{\dif V}
\newcommand{\dS}{\dif S}
\newcommand{\dA}{\dif \mathrm{A}}
\newcommand{\dmu}{\dif \mu}
\newcommand{\dnu}{\dif \nu}
\newcommand{\dtheta}{\dif \theta}
\newcommand{\domega}{\dif \omega}
\newcommand{\dsigma}{\dif \sigma}
\newcommand{\dtau}{\dif \tau}
\newcommand{\dvarphi}{\dif \varphi}
\renewcommand{\dif}{\mathrm{d}}
\renewcommand{\Dif}{\mathrm{D}}
\renewcommand{\triangle}{\Delta}
\newcommand{\ptriangle}{\partial \triangle}
\DeclareMathOperator{\im}{im}
\DeclareMathOperator{\cis}{cis}
\DeclareMathOperator{\rot}{rot}
\DeclareMathOperator{\vspan}{span}
\DeclareMathOperator{\grad}{grad}
\DeclareMathOperator{\curl}{curl}
\DeclareMathOperator{\esssup}{esssup}
\newcommand{\hodge}{{\mathnormal{\star}}}
\DeclareMathOperator{\range}{range}
\DeclareMathOperator{\Hom}{Hom}
\DeclareMathOperator{\Int}{Int}
\DeclareMathOperator{\Conv}{Conv} % convex hull
\newcommand{\ind}{\mathbbm{1}}
\DeclareMathOperator{\Res}{Res}
\DeclareMathOperator{\graph}{graph}
\DeclareMathOperator{\diam}{diam}
\DeclareMathOperator{\supp}{supp}
\DeclareMathOperator{\Vol}{Vol} % Volume
\DeclareMathOperator{\Area}{Area} % Area
\DeclareMathOperator{\area}{area}
\DeclareMathOperator{\Ind}{Ind} % Index
\newcommand{\length}{\mathrm{length}}

% logic
\DeclareMathOperator{\MOD}{MOD}
\DeclareMathOperator{\Theory}{Theory}
\newcommand{\notimplies}{%
  \mathrel{{\ooalign{\hidewidth$\not\phantom{=}$\hidewidth\cr$\implies$}}}}

% nice
\newcommand{\half}{\frac{1}{2}}
\newcommand{\pair}{\del}
\newcommand{\taking}[1]{\xrightarrow{#1}}
\newcommand{\otaking}[1]{\xleftarrow{#1}}
\newcommand{\inv}{^{-1}}
\newcommand{\ot}{\leftarrow}
\newcommand{\ninfty}{-\infty}
\newcommand{\pinfty}{+\infty}
\newcommand{\floor}[1]{\left\lfloor #1 \right\rfloor}
\newcommand{\ceil}[1]{\left\lceil #1 \right\rceil}
\newcommand{\viff}{\Big\Updownarrow} % vertical iff
\newcommand{\bmid}{\,\middle\vert\,} % big mid

% probability
\newcommand{\Prob}{\mathbf{P}}
\renewcommand{\vec}[1]{\boldsymbol{\mathbf{#1}}}
\DeclareMathOperator{\Bin}{Bin}
\DeclareMathOperator{\Geo}{Geo}
\DeclareMathOperator{\Poi}{Poi}
\DeclareMathOperator{\Exp}{Exp}
\DeclareMathOperator{\Var}{Var} % Variance
\DeclareMathOperator{\Cov}{Cov}

% special letters
\newcommand{\N}{\mathbb{N}}
\newcommand{\Z}{\mathbb{Z}}
\newcommand{\Q}{\mathbb{Q}}
\newcommand{\R}{\mathbb{R}}
\newcommand{\C}{\mathbb{C}}
\newcommand{\F}{\mathbb{F}}
\newcommand{\E}{\mathbf{E}}
\newcommand{\D}{\mathbb{D}}
\renewcommand{\H}{\mathbb{H}}
\newcommand{\ps}{\mathcal{P}}
\newcommand{\M}{\mathcal{M}}
\renewcommand{\L}{\mathcal{L}}
\renewcommand{\O}{\mathcal{O}}
\renewcommand{\U}{\mathcal{U}}
\newcommand{\Omicron}{O}
\newcommand{\powerset}{\mathcal{P}}

% text
\newcommand{\st}{\text{ s.t. }}
\newcommand{\tand}{\quad \text{and} \quad}
\newcommand{\tor}{\quad \text{or} \quad}
\newcommand{\stand}{\text{ and }}
\newcommand{\stor}{\text{ or }}
\renewcommand{\tt}[1]{\textnormal{\textbf{(#1).}}} %tt=theorem title GET RID OF

% title format
\title{\textbf{\ncourse}}

\ifx \ntype \notes
  \author{Notes taken by \nauthor\, \nemail \\ Based on lectures by \nlecturer}
  \date{\nterm\ \nyear}
\else
  \ifx \ntype \problems
    \author{These problems were solved by \nauthor \\ \nemail}
  \else
  \fi
\fi

\makeatother


\begin{document}
\maketitle

\newpage
\tableofcontents
\newpage

\section{Introduction and the Stone--Weierstrass theorem}

\begin{exercise}
  A Hamel basis for $C_{\R}([a, b])$ must be uncountable.
\end{exercise}
\begin{solution}
  One way to see this is to realize that the cardinality of the
  Hamel basis must be at least the cardinality of the biggest set of
  linearly independent vectors,
  and to see that $\set{x^r}_{r \in (0,1)}$ is an uncountable set of
  linearly independent vectors.

  Here is a more formal way of showing this.
  Assume, by way of contradiction, that $\set{v_1,v_2,\dots}$ is a Hamel
  basis for $C_{\R}([a,b])$.
  Then we have that
  \[
    X = \bigcup_{n=1}^{\infty} \vspan\set{v_1,v_2,\dots,v_n}.
  \]
  However, each set $\set{v_1,v_2,\dots,v_n}$ is nowhere
  dense in $X$ (a finite union of nowhere dense sets is nowhere dense
  and it is clear why each $\set{v_i}$ is nowhere dense).
  This implies that $X$ is a set of the first category.
  Since $C_{\R}([a,b])$ is a complete metric space,
  and $X$ is a set of the first category from Baire's theorem we
  get that $\Int(X) = \emptyset$ but this is a contradiction.
\end{solution}

\begin{exercise}
  Let $A \subseteq C_{\R}(X)$ be a subalgebra,
  and let $\overline A$ be its closure.
  Then $\overline A$ is also a subalgebra.
\end{exercise}
\begin{solution}

\end{solution}

\begin{exercise}
  Prove that the Stone--Weierstrass theorem is equivalent to
  the following statement:
  ``If $A$ is a subalgebra of $C_{\R}(X)$ which contains the constant
  functions and separates points, then $\overline A = C_{\R}(X)$''.
\end{exercise}
\begin{solution}

\end{solution}

\begin{exercise}
  Let $A \subseteq C_{\R}(X)$ be a closed subalgebra which contains the
  constant functions, and seperates points.
  For every pair of distinct points $x, y \in X$,
  and every $a, b \in \R$,
  there exists a function $g \in A$ such that $g(x) = a$ and $g(y) = b$.
\end{exercise}
\begin{solution}

\end{solution}

\begin{exercise}
  Did we use the assumption that $X$ is Hausdorff? Explain.
\end{exercise}
\begin{solution}
  
\end{solution}

\begin{exercise}
  Prove that $C(X) = \set{u +iv \mid u,v \in C_{\R}(X)}$
  where $C(X)$ is the algebra of continuous complex=valued functions on $X$.
\end{exercise}
\begin{solution}

\end{solution}

\begin{exercise}
  Is the space $C[x,y]$ of two variable complex polynomials dense in
  $C(\overline \D)$.
\end{exercise}
\begin{solution}

\end{solution}

\begin{exercise}
  Let $f \in C_{\R}([-1, 1])$ such that $f(0) = 0$.
  Prove that $f$ can be approximated in uniform norm by
  polynomials with zero constant coefficient.
\end{exercise}
\begin{solution}

\end{solution}Prove that if A is a closed subalgebra of CH(X)
that separates points and contains all the constant H-valued functions, then A = CH(X).

\begin{exercise}
  Let $C_0(\R)$ be the algebra of all continuous functions that vanish
  at infinity, that is
  \[
    C_0(\R) = \set{f \in C(\R) \mid
        \lim_{x \to \infty} f(x) = \lim_{x \to \ninfty} f(x) = 0}.
  \]
  Let $A$ be a self-adjoint subalgebra of $C_0(\R)$ that separates points,
  and assume that for all $x \in \R$,
  there exists $f \in A$ such that $f(x) \neq 0$.
  Prove that $A = C_0(\R)$.
\end{exercise}
\begin{solution}

\end{solution}

\begin{exercise}
  Determine which of the following statements are true.
  \begin{enumerate}
  \item[(1)] For every $f \in C_{\R}([a, b])$,
    all $x_1, x_2, \dots, x_n \in [a,b]$,
    and every $\epsilon > 0$,
    there exists a polynomial $p$ such that
    $\sup_{x \in [a,b]} |p(x)-f(x)| < \epsilon$
    and $p(x_i) = f(x_i)$ for $i = 1,2,\dots,n$.
    \item[(2)] For every $f \in C_{\R}([a, b])$
    which has $n$ continuous derivatives on $[a,b]$,
    and every $\epsilon > 0$,
    there exists a polynomial $p$ such that
    $\sup_{x \in [a,b]} |p(x)-f(x)| < \epsilon$
    and $p^{(k)}(x_i) = f^{(k)}(x_i)$ for $k = 0,1,\dots,n$.
  \end{enumerate}
  Determine when happens in the case we replace the word ``polynomial'' with
  ``trigonometric polynomial''.
\end{exercise}
\begin{solution}

\end{solution}

\begin{exercise} \phantom{}
  \begin{enumerate}
    \item[(1)] Let $f$ be a continuous real-valued function on $[a,b]$
      such that $\int_{a}^{b} f(x) x^n \dx = 0$
      for all $n = 0,1,2,\dots$.
      Prove that $f \equiv 0$.
    \item[(2)] Let $f$ be a continuous real-valued function on $[-1,1]$
      such that $\int_{-1}^{1} f(x) x^n \dx = 0$
      for all odd $n$. What can you say about $f$?
  \end{enumerate}
\end{exercise}
\begin{solution}

\end{solution}

\begin{exercise}
  A $\Z^k$-periodic function is a function $f \colon \R^k \to \R$
  such that $f(x + n) = f(x)$ for all $x \in \R^k$ and $n \in \Z^k$.
  A trigonometric polynomial on $\R^k$ is a function $q$ of the form
  \[
    q(x) = \sum_{n} a_n \cos(2\pi n \cdot x) + b_n \sin(2\pi n \cdot x)
  \]
  where the sum is finite and $n \cdot x$ is the scalar product
  of $x \in \R^k$ and $n \in \Z^k$.
  Prove the trigonometric approximation theorem
  in the multidimensional case.
\end{exercise}
\begin{solution}

\end{solution}

\begin{exercise}[Holladay]
  Let $X$ be a compact Hausdorff space and let $C_{\H}(X)$ denote the
  algebra of continuous functions taking quaternionic values.
  Here, $\H$ denotes the quaternions, that is,
  the four-dimensional real vector space
  \[
    \H = \set{a + bi + cj + dk \mid a,b,c,d \in \R}
  \]
  with a basis consisting of $1 \in \R$
  and three additional elements $i$, $j$, $k$, and with multiplication
  determined by the multiplication of the reals and by the identities
  \[
    i^2 = j^2 = k^2 = ijk = -1.
  \]
  We let the topology of $\H$ to be the one given by identifying it with $\R^4$.
  Prove that if $A$ is a closed subalgebra of $C_{\H}(X)$
  that separates points and contains all the constant $\H$-valued functions,
  then $A = C_{\H}(X)$.
\end{exercise}
\begin{solution}

\end{solution}

\newpage

\section{Hilbert spaces}

\begin{exercise}
  Prove the following polariztion identities.
  \begin{enumerate}
    \item[(1)] If $G$ is a real inner product space, then for all $f,g \in G$,
      \[
        \ip{f}{g} = \frac{1}{4} \del{\norm{f + g}^2 - \norm{f - g}^2}.
      \]
    \item[(1)] If $G$ is a complex inner product space,
      then for all $f,g \in G$,
      \[
        \ip{f}{g} = \frac{1}{4} \del{\norm{f + g}^2 - \norm{f - g}^2
          + i \norm{f + ig}^2 - i \norm{f - ig}^2}.
      \]
  \end{enumerate}
\end{exercise}
\begin{solution}

\end{solution}

\begin{exercise}
  Let $G$ be an inner product space, let $g \in G$,
  and let $F$ be a dense subset of $G$.
  Show that if $\ip{f}{g} = 0$ for all $f \in F$, then $g = 0$.
\end{exercise}
\begin{solution}

\end{solution}

\begin{exercise}
  Show that for every $f \in PC[a, b]$,
  there is at most one $g \in C([a, b])$ in its equivalence class.
\end{exercise}
\begin{solution}

\end{solution}

\begin{exercise}
  Show that with respect to the inner product
  \[
    \ip{f}{g} = \int_{a}^{b} f(t) \overline{g(t)} \dt
  \]
  that $PC[a, b]$ is not complete.
\end{exercise}
\begin{solution}

\end{solution}

\begin{exercise}
  Prove that a function in $C([a,b])$ is completely determined by the values of
  its integrals over sub-intervals of $[a,b]$.
\end{exercise}
\begin{solution}

\end{solution}

\begin{exercise}
  Show that a function in $L^2([a,b])$
  is completely determined by the values of
  its integrals over sub-intervals of $[a,b]$.
\end{exercise}
\begin{solution}

\end{solution}

\begin{exercise}[tricky]
  Try to justify the statement: ``for every $r < 1/2$,
  the function given by $f(x) = x^{-r}$ for $x \in (0, 1]$
  (and undefined for $x = 0$) is in $L^2[0,1]$.''
\end{exercise}
\begin{solution}

\end{solution}

% additional exercises

\newpage

\section{Orthogonality, projections, and bases}

\begin{exercise}
  Prove that for any set $S$, $S^{\perp}$ is a closed subspace,
  and that $S \cap S^{\perp} \subseteq \set{0}$.
\end{exercise}
\begin{solution}

\end{solution}

\begin{exercise}
  Prove that for any set $S$, $S^{\perp}$ is a closed subspace,
  and that $S \cap S^{\perp} \subseteq \set{0}$.
\end{exercise}
\begin{solution}

\end{solution}

\begin{exercise}
  Prove that a finite dimensional subspace of an inner product space is
  closed.
\end{exercise}
\begin{solution}

\end{solution}

\begin{exercise}
  Prove that in an inner product space $G$,
  a series $\sum_{i \in I} v_i$ converges to $g \in G$ if
  and only if there exists a countable set $J$ contained in $I$ such that
  \begin{enumerate}
    \item[(a)] $v_i = 0$ if $i \in I \setminus J$;
    \item[(b)] for any enumeration $J = \set{j_1,j_2,\dots}$ of the elements
      of $J$, the limit $\lim_{N \to \infty} \sum_{n=1}^{N} v_{j_n} = g$
      exists in $G$.
  \end{enumerate}
\end{exercise}
\begin{solution}

\end{solution}

\begin{exercise}
  Suppose that $a_i \geq 0$ for all $i$.
  Prove that $\sum_{i \in I} a_i$ converges and is bounded
  by $M$ if and only if the set of all finite sums
  \[
    \set{\sum_{i \in F} a_i \colon F \subset I \text{ finite}}
  \]
  is bounded by $M$.
\end{exercise}
\begin{solution}

\end{solution}

\begin{exercise}
  Deduce the Cauchy--Schwarz inequality from Bessel's inequality
  (this should be one line).
  Did this involve circular reasoning?
\end{exercise}
\begin{solution}

\end{solution}

% idk

\begin{exercise}
  Two Hilbert spaces are isomorphic if and only if they have the same di-
mension.
\end{exercise}
\begin{solution}

\end{solution}

\begin{exercise}
  Prove that a Hilbert space is separable if and only if its dimension is less
than or equal to $\alpha_0$
(recall that a metric space is said to be separable if it contains a countable
dense subset).
\end{exercise}
\begin{solution}

\end{solution}

\begin{exercise}
  Prove that any Hilbert space is isomorphic to $\ell^2(I)$ for some set $I$.
\end{exercise}
\begin{solution}

\end{solution}

\begin{exercise}
  Let $G$ be the linear span of all functinos of the form
  $\sum_{k=1}^{n} a_k e^{i \lambda_k t}$, where $\lambda_k \in \R$.
  On $G$, we define a form
  \[
    \ip{f}{g} =
    \lim_{T \to \infty} \frac{1}{2T} \int_{-T}^{T} f(t) \overline{g(t)} \dt.
  \]
  Prove that $G$ is an inner product space.
  Let $H$ be the completion of $G$.
  Prove that H is not separable.
  Find an orthonormal basis for $H$.
\end{exercise}
\begin{solution}

\end{solution}


\end{document}
